\documentclass[x11names]{beamer}
\usepackage{etex,beamerfoils}
\usepackage{array,pifont,amsmath,latexsym,epsfig,amsfonts,yfonts}
\usepackage{amstext,amsopn,amsxtra,amsgen,amsbsy,amscd,amssymb,dcolumn,graphicx,setspace,pgfplots,booktabs,natbib}
\usepackage[normalem]{ulem}
\usepackage[all]{xy}
\beamertemplatetheoremsnumbered
\setbeamertemplate{caption}[numbered]
\usetheme{Copenhagen}
\definecolor{NewBlue}{RGB}{4, 114, 155}
\definecolor{magenta}{cmyk}{0,1,0,0}
\usecolortheme[named=NewBlue]{structure}
%\usecolortheme[named=SkyBlue4]{structure}

\newcommand{\D}{\displaystyle}
\newcommand{\E}{\begin{eqnarray*}}
\newcommand{\F}{\end{eqnarray*}}
\newcommand{\EE}{\begin{eqnarray}}
\newcommand{\FF}{\end{eqnarray}}
\newcommand{\IZ}{\begin{itemize}}
\newcommand{\ZI}{\end{itemize}}
\newcommand{\EN}{\begin{enumerate}}
\newcommand{\NE}{\end{enumerate}}
\newcommand{\itemc}{\item[$\circ$]}
\newcommand{\IZdash}{\begin{itemize} \renewcommand{\labelitemi}{-}}
\newcommand{\tn}{\textnormal}

% New commands to assist Beamer
\newcommand{\itemi}{\item<1->}
\newcommand{\itemii}{\item<2->}
\newcommand{\itemiii}{\item<3->}
\newcommand{\itemiv}{\item<4->}
\newcommand{\itemv}{\item<5->}
\newcommand{\itemvi}{\item<6->}
\newcommand{\itemvii}{\item<7->}
\newcommand{\itemviii}{\item<8->}
\newcommand{\itemix}{\item<9->}
\newcommand{\itemx}{\item<10->}

\newcommand{\itemci}{\item[$\circ$]<1->}
\newcommand{\itemcii}{\item[$\circ$]<2->}
\newcommand{\itemciii}{\item[$\circ$]<3->}
\newcommand{\itemciv}{\item[$\circ$]<4->}
\newcommand{\itemcv}{\item[$\circ$]<5->}
\newcommand{\itemcvi}{\item[$\circ$]<6->}
\newcommand{\itemcvii}{\item[$\circ$]<7->}
\newcommand{\itemcviii}{\item[$\circ$]<8->}
\newcommand{\itemcix}{\item[$\circ$]<9->}
\newcommand{\itemcx}{\item[$\circ$]<10->}


\newcommand{\fr}{\frame}
\newcommand{\ft}{\frametitle}
\newcommand{\ons}{\onslide}
\newcommand{\tc}{\textcolor}

\pgfplotsset{height =7cm,compat=1.7,width=7cm,
    standard/.style={
    	unit vector ratio*=1 1 1,
        axis x line=middle,
        axis y line=middle,
        enlarge x limits=0.15,
        enlarge y limits=0.15,
        every axis x label/.style={at={(current axis.right of origin)},anchor=north west},
        every axis y label/.style={at={(current axis.above origin)},anchor=north east}
    }
}
\newcommand{\drawge}{-- (rel axis cs:1,0) -- (rel axis cs:1,1) -- (rel axis cs:0,1) \closedcycle}
\newcommand{\drawle}{-- (rel axis cs:1,1) -- (rel axis cs:1,0) -- (rel axis cs:0,0) \closedcycle}

\usetikzlibrary{patterns}
\usetikzlibrary{arrows}
\usetikzlibrary{shapes}

\renewcommand{\bibsection}{}

% Watermark example: https://texblog.org/2012/02/17/watermarks-draft-review-approved-confidential/
\usepackage{draftwatermark}
\SetWatermarkText{This content is protected and may not be shared uploaded or distributed.}
\SetWatermarkScale{0.25}
\SetWatermarkAngle{0}
\SetWatermarkColor[gray]{1.0}
\setbeamercolor{background canvas}{bg=}%transparent canvas

\title[New-Keynesian Modeling\hspace{2em}\insertframenumber/
\inserttotalframenumber]{Introduction to New-Keynesian Business Cycle Modeling}
\author{Brian C.~Jenkins \\ ~ \\University of California, Irvine}

\begin{document}

\frame{\maketitle}

\fr{
\ft{Introduction}

\IZ
\item We analyze a model that is representative of the \textbf{new-Keynesian} class of business cycle models.

\

\itemii The model is \emph{Keynesian} because it assumes that nominal prices are not flexible; they adjust gradually in response to an exogenous shock to the economy.

\

\itemiii The model is \emph{new} because it is built upon a solid microeconomic foundation.

\

\itemiv New-Keynesian models are useful for explaining the equilibrium relationships between inflation, output, and interest rates.

\ZI
}


\fr{
\ft{The Euler Equation}

\IZ
\item Suppose that a \emph{representative household} lives for two periods.

\

\itemii The household receives utility from consuming goods in each period.

\

\itemiii The \emph{lifetime utility} to the household from consuming $C_0$ in the first period and $C_1$ in the second is denoted by $U(C_0,C_1)$ and is written as:
	\EE
	U(C_0,C_1) & = & u(C_0) + \beta u(C_1),  \label{eqn:two_per_utility}
	\FF
	
where $u(\cdot)$ is the period utility function with $u'(\cdot)>0$, $u''(\cdot)<0$ and $\beta \leqslant 1$.

\ZI
}


\fr{
\input{figure_015_intro_nk_modeling_flow_utility.tex}
}

\fr{
\begin{figure}[t] \caption{\label{fig:two_per_flow_utility} \textbf{The household's indifference curves} over combinations of consumption in periods 0 and 1.}


	\begin{center}

	\begin{tikzpicture}
	\begin{axis}[standard,%tiny,
    %scale only axis,
    font=\small,
    ytick={0},
    yticklabels={0},
    xtick={0},
    xticklabels={0},
    minor tick num=0,
    axis y line=left,
    axis x line=middle,
    xmin = 0,
    xmax = 5,
    ymin = 0,
    ymax = 5,
    xlabel style={at=(current axis.right of origin), anchor=west,text width=0 cm},
    ylabel style={left},
    xlabel=$C_0$,ylabel=$C_1$,
%    title={(B) Indifference curves}
	]
	%\addplot[smooth,blue,mark=none,line width = 1.5pt,
	%domain=0:4.25,samples=100]
	%{exp(0.25-0.95*ln(x))} node [pos=1,inner sep=2pt,label=right:$\bar{u}_0$] {};
	\addplot[smooth,blue,mark=none,line width = 1.5pt,
	domain=0:4.25,samples=100]
	{exp(.5-0.95*ln(x))} node [pos=1,inner sep=2pt,label=right:$\bar{U}_1$] {};
	\addplot[smooth,blue,mark=none,line width = 1.5pt,
	domain=0:4.25,samples=100]
	{exp(1.25-0.95*ln(x))} node [pos=1,inner sep=2pt,label=right:$\bar{U}_2$] {};
	\addplot[smooth,blue,mark=none,line width = 1.5pt,
	domain=0:4.25,samples=100]
	{exp(1.775-0.95*ln(x))} node [pos=1,inner sep=2pt,label=right:$\bar{U}_3$] {};
	%\addplot[smooth,blue,mark=none,line width = 1.5pt,
	%domain=0:4.25,samples=100]
	%{exp(-0.95*ln(x))} node [pos=1,inner sep=2pt,label=right:$\bar{u}_0$] {};
	
	%\addplot[smooth,blue,mark=none,line width = 1.5pt,
	%domain=1:4,samples=100]
	%{exp(ln(2)+ln(2.5)-ln(x))} node [pos=0.95,inner sep=2pt,label=right:$\bar{U}^*$] {};
	%\addplot[blue,dotted,mark=,line width = 1pt]
	%coordinates{(0,3) (3,3) (3,0)};
	%\addplot[blue,solid,mark=*,line width = 1pt]
	%coordinates{(3,3)} node[pos=0.5, above, font=\footnotesize, yshift=0pt,xshift = 25pt] {endowment};
	%\addplot[magenta,dotted,mark=,line width = 1pt]
	%coordinates{(2,0) (2,2.5) (0,2.5)};
	%\addplot[magenta,solid,mark=*,line width = 1pt]
	%coordinates{(2,2.5)} ;
	\addplot[->,black,mark=none,line width = 1pt]
	coordinates{(1,1) (3,3)} node [pos=1,inner sep=2pt,xshift=5pt,label=above:increasing $U$] {};
	\end{axis}
	\end{tikzpicture}
	\end{center}
\end{figure}
}

\fr{
\ft{The Euler Equation}

\IZ
\item Household has no initial wealth

\

\itemii Receives an income \emph{endowment} of $Y_0$ in period 0 and $Y_1$ in period 1.

\

\itemiii In period 0, household chooses an amount of $Y_0$ to save $S_1$ given a real interest rate $r_0$.
\ZI
}


\fr{
\ft{The Euler Equation}

\IZ
\item The household's budget constraints for periods 0 and 1 are:
	\begin{align}
	C_0 & = Y_0 - S_1,  \label{eqn:two_per_bc0}\\
	C_1 & = Y_1 + (1+r_0)S_1.  \label{eqn:two_per_bc1}
	\end{align}
	
\

\itemii Combine the two period budget constraints to obtain the intertemporal budget constraint:

	\EE
	C_1 & = & Y_1 + (1+r_0)Y_0 - (1+r_0)C_0.  \label{eqn:two_per_intertemporal_bc}
	\FF

\ZI
}

\fr{
\ft{The Euler Equation}


\IZ
\item The household's optimization problem is:
	\EE
	&&\max_{C_0,C_1,S_1} \; u(C_0) + \beta u(C_1)\\
	&& \; \; \; \;  \; \; \; \; \tn{s.t.} \; \; \; \; C_0 = Y_0 - S_1\\
	&&\; \; \; \; \; \; \; \; \; \; \; \; \; \; \; \; \; C_1 = Y_1 + (1+r_0)S_1.
	\FF

\

\itemii Equivalent to:
	\EE
	\max_{S_1} \; u( Y_0 - S_1) + \beta u\left[ Y_1 + (1+r_0)S_1\right]
	\FF
\ZI

}

\fr{
\ft{The Euler Equation}


\IZ
\item Setting derivative of the objective function with with respect to $S_1$ equal to zero:
	\EE
	-u'( Y_0 - S_1) + (1+r_0) \beta u'\left[ Y_1 + (1+r_0)S_1\right] & = & 0 \label{eqn:two_per_s1_foc}
	\FF
where $u'(\cdot)$ denotes the period marginal utility of consumption. 


\

\itemii Equation (\ref{eqn:two_per_s1_foc}) is the \emph{first-order condition} for the optimal choice of $S_1$.

\

\itemiii We can rewrite the first-order condition for $S_1$ in terms of consumption:
	\EE
	\boxed{u'(C_0) = \beta (1+ r_0) u'(C_1)} \label{two_per_euler}
	\FF


\ZI
}

\fr{
\ft{The Euler Equation}


\IZ
\item Given $Y_0$, $Y_1$, and $r_0$, the Euler equation (\ref{two_per_euler}) determines whether the household will borrow or save.

\

\item If $Y_0$ is relatively large or if $r_0$ is relatively large, then the household will want to save.
\ZI
}


\fr{
\begin{figure}[t] \caption{\label{fig:two_per_opt_saving} \textbf{A Household that saves}. The period 0 endowment is large relative to the period 1 endowment}


	\begin{center}

	\begin{tikzpicture}
	\begin{axis}[standard,%tiny,
    %scale only axis,
    font=\small,
    ytick={0,1,2},
    yticklabels={0,$Y_1$,$C_1^*$},
    xtick={0,2,3},
    xticklabels={0,$C_0^*$,$Y_0$},
    minor tick num=0,
    axis y line=left,
    axis x line=middle,
    xmin = 0,
    xmax = 5,
    ymin = 0,
    ymax = 5,
    xlabel style={at=(current axis.right of origin), anchor=west,text width=0 cm},
    ylabel style={left},
    xlabel=$C_0$,ylabel=$C_1$,
%    title={A Household that saves}
	]
	\addplot[smooth,blue,mark=none,line width = 1.5pt,
	domain=0:4.25,samples=100]
	{4-x} node [pos=1,inner sep=2pt,label=right:$\bar{u}_0$] {};
	\addplot[smooth,blue,mark=none,line width = 1.5pt,
	domain=0:4.25,samples=100]
	{exp(ln(2) + ln(2) - ln(x))} node [pos=1,inner sep=2pt,label=right:$\bar{U}^*$] {};
	\addplot[black,dotted,mark=,line width = 1pt]
	coordinates{(0,1) (3,1) (3,0)};
	\addplot[black,solid,mark=*,line width = 1pt]
	coordinates{(3,1)} node[pos=0.5, below, font=\footnotesize, yshift=0pt,xshift = -5pt] {e};
	\addplot[magenta,dotted,mark=,line width = 1pt]
	coordinates{(2,0) (2,2) (0,2)};
	\addplot[magenta,solid,mark=*,line width = 1pt]
	coordinates{(2,2)} ;
	\end{axis}
	\end{tikzpicture}
	\end{center}
\end{figure}
}

\fr{
\input{figure_015_intro_nk_modeling_borrowing.tex}
}


\fr{
\ft{The Euler Equation}


\IZ
\item We will assume that the period utility function is the natural logarithm of consumption:
	\EE
	u(C) & = & \log C.
	\FF
	
\

\itemii With this utility function, the household's Euler equation is expressed as:
	\EE
	\boxed{\frac{1}{C_0} = \frac{\beta (1+ r_0)}{C_1}} \label{two_per_euler}
	\FF
\ZI
}



\fr{
\ft{The Demand for Goods}

\IZ
\item Assume no investment, government purchases, or net exports, so $Y = C$ in every period:
	\EE
	\frac{1}{Y_0} & = & \beta \frac{1}{Y_1} (1+r_{0})
	\FF

\

\itemii Next, take the log of each side to obtain a linear representation of the Euler equation:
	\EE
	-\log Y_0 & = & \log \beta - \log Y_1 + \log \left(1+r_{0}\right)
	\FF
\ZI	

}


\fr{
\ft{The Demand for Goods}

\IZ
\itemi Rearranging and defining $\bar{r} \equiv -\log \beta$ as the natural rate of interest and $y \equiv \log Y$, obtain:
	\EE
	y_0 & = & y_1 - (r_{0}-\bar{r}).
	\FF

\itemii Next:
	\IZ
	\itemiii The previous expression will hold for any two adjacent time periods $t$ and $t+1$
	\itemiv Incorporate uncertainty: replace $y_1$ with its expectation conditional on period 0 information $E_0y_1$
	\itemv Append an exogenous demand shock
	\ZI

\itemvi Finally:
	\EE
	\boxed{y_t = E_ty_{t+1} -(r_{t} -\bar{r}) + g_t}
	\FF

\ZI
}

\fr{
\ft{Equilibrium}

\IZ
\item The complete set of equilibrium conditions:
\ZI

\EE
y_t & = & E_t y_{t+1}  - \left( r_{t} - \bar{r}\right) + g_t\\
\pause i_{t} & = & r_{t} + E_t \pi_{t+1}\\
\pause i_{t} & = & \bar{r} + \pi^T + \phi_{\pi}\big(\pi_t - \pi^T\big) + \phi_{y}\big(y_t - \bar{y}\big) + v_t\\
\pause \pi_t -\pi^T & = & \beta \left( E_t\pi_{t+1} - \pi^T\right)  + \kappa (y_t -\bar{y})+ u_t,
\FF

\IZ
\itemv The equations are: dynamic IS equation, Fisher equation, monetary policy rule, new-Keynesian Phillips curve or dynamic AS equation.
\ZI
}





%\fr{
%\ft{References}
%\bibliographystyle{aer}
%\bibliography{comp_econ_references.bib}
%}

\end{document}